\section{后端关系、支持编码与构造费用}
\subsection{四状态监视器的完整定义}\label{app:automaton}
式~\eqref{eq:language} 的部分 DFA 状态为 $0,1,2,3$，初态为 $0$，接受态为 $1,2$。非空转移如下，未列转移均不存在：
\begin{align*}
 s:&\quad 0\to1,\ 1\to1,\\
 r:&\quad 0\to1,\ 1\to1,\ 2\to2,\\
 g:&\quad 1\to2,\ 2\to2,\ 3\to2,\\
 \bar r:&\quad 1\to3,\ 2\to3.
\end{align*}
状态 $1$ 记录已经读过非空 $s/r$ 前缀，仍可延长该前缀；读到 $g$ 后进入只能读一致性后缀的状态 $2$；$\bar r$ 进入状态 $3$ 并强制随后至少一个 $g$。后缀中的 $r$ 可以继续停留状态 $1$，因为下一条 $s$ 仍可归入更长的因果前缀；一旦出现 $g$ 或完成 $\bar r g^+$，就不能再读 $s$。由此该 DFA 恰识别式~\eqref{eq:language}。

对每个事件 $e$，安装标记 $(e,1)\to(e,0)$ 和 $(e,2)\to(e,0)$。若事件图存在接受语言中的闭合游走，从 $(e,0)$ 沿乘积普通边到 $(e,1)$ 或 $(e,2)$，接一条标记便闭合。反之，任何恰一标记闭合游走可以旋转到该标记之后开始，其余全部普通边给出从初态到接受态、起止事件相同的游走。故二者等价。CO 监视器只有一个状态，每条 $g$ 边原样加入并要求普通图无环。

必须计数恰好一个标记。例如普通边 $0\to3,2\to1$，标记边 $1\to0,3\to2$ 形成两标记环，却没有恰一标记闭合游走。把普通边与标记边混合取闭包，或者允许 $T;T$，会错误拒绝这个关系状态。

\subsection{增量更新及端口投影}\label{app:update}
以下 $P$ 含普通恒等路径，$T$ 恰含一标记路径。插入普通边 $u\to v$ 前，若 $T(v,u)$ 则出现一标记环，应拒绝；否则用插入前的关系计算
\begin{align}
 P'&=P\cup P(*,u)\times P(v,*),\label{eq:updateP}\\
 T'&=T\cup T(*,u)\times P(v,*)\cup P(*,u)\times T(v,*).\label{eq:updateT}
\end{align}
普通路径中的新边若重复出现，可删去其间普通回路而不改变可达性。恰一标记路径若两次使用新边且标记处于其间，则存在旧 $T(v,u)$，已被拒绝；否则重复普通回路同样可去。于是列出的两类含标记扩展已经穷尽。

插入标记 $u\to v$ 前，若 $P(v,u)$ 则拒绝；否则
\[
 T'=T\cup P(*,u)\times P(v,*),\qquad P'=P.
\]
新标记不能与旧 $T$ 路径串接。CO 单状态图插普通边时检查 $P(v,u)$，随后使用式~\eqref{eq:updateP}。

每组完成后，投影到下一层由全部候选决定的 $I\times O$。装载下一层时恢复这些条目和普通恒等关系，再加入当前组新端点。可拼接的旧矩形条目仍在同一矩形中，其普通／一标记闭包已由前一层保存；未保留端点不可能再成为未来边的接口。引理~\ref{lem:ports} 的新旧段分解因此适用于投影后的更新。整个状态键还包括所有监视器和活跃读计数；逐监视器独立挑选代表历史不具有这个性质。

\subsection{支持 CNF 的完整构造}\label{app:supportcnf}
令接受 DAG 有唯一起点 $s_0$ 和终点 $t_0$，边 $e$ 的原始文字标签为 $\ell(e)$。固定文字可以作为公共起始边；每个原始变量在每条完整路径上恰被赋值一次。定义阻断条件
\[
 D_e=\bigvee_{\lambda\in\ell(e)}\neg\lambda.
\]
空标签使 $D_e$ 恒假。引入前缀不可达位 $A_v$ 与后缀不可达位 $Z_v$，按 DAG 顺序定义
\begin{align*}
 A_{s_0}&=\bot,& A_v&=\bigwedge_{e:u\to v}(A_u\lor D_e),\quad v\ne s_0,\\
 Z_{t_0}&=\bot,& Z_v&=\bigwedge_{e:v\to w}(D_e\lor Z_w),\quad v\ne t_0.
\end{align*}
空合取为真。对每个原始文字定义无支持位
\[
 W_\lambda=\bigwedge_{e:\lambda\in\ell(e)}(A_{\mathrm{src}(e)}\lor D_e\lor Z_{\mathrm{dst}(e)}).
\]
把这些无环电路以双向 Tseitin 子句编码，再加 $\neg Z_{s_0}$ 和每个 $W_\lambda\Rightarrow\neg\lambda$。常量的相反文字若没有任何边携带，其 $W$ 为空合取，会正确传播对应常量。电路和 CNF 大小均线性于图与标签出现次数。

在部分原始赋值 $\alpha$ 下，所有被阻断边的 $D_e$ 都由已知反文字传播为真。沿拓扑序归纳：若某节点无相容前缀，则所有入边要么被阻断，要么来自已传播不可达的前驱，故 $A_v$ 被传播为真。反向归纳得到所有无相容后缀的 $Z_v$。无接受路径时 $Z_{s_0}$ 为真，与单位子句冲突。若 $\lambda$ 不支持，任何携带它的边必在前缀、边阻断、后缀三种条件之一失败，因而 $W_\lambda$ 传播为真并推出 $\neg\lambda$。这个推导只用真的单调传播，不把未赋值文字当作假。

完整赋值 $x$ 若没有接受路径，会发生上述冲突；若有接受路径，则按电路的真实值赋辅助变量。任何 $x$ 为真的文字都在这条路径中出现，因此其 $W$ 为假，其余 $W_\lambda\Rightarrow\neg\lambda$ 也满足，得到 CNF 模型。这再次证明全赋值投影。由该精确投影，传播不会排除合法完成；支持不可能的反文字被排除，恰给出全部必然原始文字。该强度是相对原始变量的域一致性，与允许任意赋值辅助变量的传播完备性有区别~\cite{kucera2021}。

\subsection{协议构造的保守位复杂度}\label{app:complexity}
令 $L_{\rm in}$ 为包含事件、二进制值与容量、元数据和显式接口描述的输入位长度，$\ell=L_{\rm in}+|X|+1$。在二进制随机访问模型下，可用如下朴素实现得到 $O(\ell^8)$ 的预处理及复核上界，证书和完整计划编码长度为 $O(\ell^6)$。这是保守构造界，并非实验实现的实测增长拟合。

注册表枚举至多 $O(n^2)$ 个读写／写写对。两个协议的来源下界和候选上界只对显式事件实例化；每条候选可以列出至多 $n$ 个前提来源，代价至多 $O(n^3)$ 次关系或数值检查。一个容量值即使很大，算法也不遍历 $0,1,\ldots,C$；必要的数组地址及初始化必须已经显式列入输入。整数比较、加减及取模保守以两个 $O(\ell)$ 位整数的 $O(\ell^2)$ 位操作计费。

永久来源域只减，$H/L$ 只增，全部可能变化不超过 $O(n^2)$ 项。每次有变化的遍历用 $O(n^3)$ 次单元操作做闭包及全量 R1/R2/$R_\forall$ 检查，故饱和至多 $O(n^5)$ 次单元操作，连同数值和引用操作由 $O(\ell^7)$ 覆盖。地址闭合分区满足 $\sum_i|E_i|\le n$，局部及全局饱和总费用仍在同阶粗界中。覆盖边计算、前沿区间扫描、布局拓扑序、失败串登记和计划检查均为更低阶的显式多项式过程。

推导记录对每个永久事实只保存第一次理由，共 $O(n^2)$ 条。R1/R2 与闭包记录常数个先前事实，$R_\forall$ 最多列出 $n$ 个当前候选及共同前序；域用删除日志和版本标识重放，也可保守保存每轮矩阵快照。规范计划至多 $O(n^2)$ 层，每层列出至多 $n$ 个端点及归属引用，其长度为 $O(n^3\log(\ell+2))$。这些材料由 $O(\ell^6)$ 位覆盖；再独立重放记录及计划，保守 $O(\ell^8)$ 足够。复核检查的是已证明规则的有限前提，不调用“对所有完成成立”的语义 oracle。

接上式~\eqref{eq:size} 的后端，在固定协议参数 $\vartheta$ 下，只要相应协议界给出 $K\le g(\vartheta)$，总构造可写为 $f(\vartheta)\ell^{O(1)}$。参数分别包含缓冲的容量、地址数、工作量及在阶段版本中使用的失败次数界，或纯失败链版本的容量、地址数、工作量；屏障包含参与线程数、地址数、工作量；组合还包含组件数及每个实际前沿参数。这个结论不把协议识别、任意布局优化或未支持守卫的处理视为免费，也不把无界失败链的袋界代入需要原始前沿界的组合入口。

\subsection{原始 SAT 核对的投影理由}\label{app:oracle}
独立原始编码采用 RF 恰一、逐地址严格总 CO、输入守卫及一个辅助关系 $H_o$。令条件基本边 $B=\SB_{\rm adj}\cup\rf\cup\init$，加入 $B(u,v)\Rightarrow H_o(u,v)$、$H_o(u,v)\land B(v,w)\Rightarrow H_o(u,w)$ 及 $\neg H_o(u,u)$。同时用 WW／WR／RW／RR 子句约束 $H_o$ 与原始 RF／CO。它没有导入协议证书的来源删域或共同关系。

固定原始赋值后，按路径长度归纳有 $B^+\subseteq H_o$。四类一致性子句都以 $H_o$ 为前提，故把它缩小为 $B^+$ 不破坏子句；反自环保证实际 HB 无环，从而得到原始 RA 完成。反向，合法完成取 $H_o=B^+$ 就满足所有子句。因此编码的存在辅助变量投影是原始语义。这个证明解释了负支持核查为何不循环依赖待测证书；具体测量仍是选定样本的求解结果，不等同于对所有输入机械证明实现。
